\begin{figure}[t]
\centering
\resizebox{\linewidth}{!}{%
\begin{tikzpicture}[
  node distance=1.45cm and 1.9cm,
  block/.style={draw, rounded corners, align=center, minimum width=3.0cm, minimum height=0.95cm, fill=gray!8},
  knowledge/.style={draw, rounded corners, align=center, minimum width=3.0cm, minimum height=0.9cm, fill=blue!7},
  model/.style={draw, rounded corners, align=center, minimum width=3.0cm, minimum height=0.95cm, fill=green!8},
  output/.style={draw, rounded corners, align=center, minimum width=3.0cm, minimum height=0.95cm, fill=orange!10},
  arrow/.style={-{Latex[length=2mm]}, thick}
]
\node[block] (q) {Medical question\\$q$};
\node[block, right=of q] (first) {First-stage retrieval\\BM25 + dense + RRF};
\node[block, right=of first] (pool) {Top-$M$ candidate\\evidence passages};
\node[knowledge, below=of first] (kg) {Medical constraints\\MeSH hierarchy\\entities\\PrimeKG auxiliary relations};
\node[knowledge, below=of pool] (hg) {Local hypergraph\\question, passage, document,\\entity, concept nodes};
\node[model, right=of pool] (ltr) {Query-group\\LambdaMART reranker};
\node[output, right=of ltr] (out) {Ranked evidence\\top-$k$ passages\\optional paths};
\draw[arrow] (q) -- (first);
\draw[arrow] (first) -- (pool);
\draw[arrow] (pool) -- (ltr);
\draw[arrow] (ltr) -- (out);
\draw[arrow] (kg) -- (hg);
\draw[arrow] (pool) -- (hg);
\draw[arrow] (q) |- (hg);
\draw[arrow] (hg) -- node[below, align=center] {diffusion, centrality,\\hierarchy features} (ltr);
\end{tikzpicture}}
\caption{Overview of KCH-MedRank. The method first builds a high-recall candidate pool, constructs a local knowledge-constrained hypergraph around each question and candidate set, and uses hypergraph and biomedical semantic features in query-group learning-to-rank.}
\label{fig:method-overview}
\end{figure}
